\chapter{Resultados y discusion}
\label{chap:resultados}

\section{Hallazgo principal}

El hallazgo principal es que los modelos no lineales superaron de forma clara a la regresion lineal en la estimacion de volatilidad implicita fuera de muestra. En las ejecuciones finales disponibles, Random Forest obtuvo el mejor resultado de test, con MAE 0,054366, RMSE 0,088403 y $R^2$ 0,790941. XGBoost quedo cerca, con RMSE 0,091735 y $R^2$ 0,774885. Las redes inspiradas en tensor-train mejoraron al baseline lineal, aunque quedaron por debajo de los ensembles de arboles en la mejor configuracion observada.

\section{Resultados cuantitativos}

La tabla \ref{tab:resultados-modelos} resume las metricas finales. Los valores proceden de los metadatos de reentrenamiento final sobre trainval y evaluacion en test.

\begin{table}[h]
\centering
\begin{tabular}{lrrr}
\toprule
Modelo & MAE test & RMSE test & $R^2$ test \\
\midrule
Regresion lineal & 0,067585 & 0,183824 & 0,096059 \\
Regresion lineal progresiva & 0,070544 & 0,187747 & 0,057069 \\
Red tensor-train & 0,056056 & 0,143025 & 0,452787 \\
Red tensor-train progresiva & 0,055597 & 0,117404 & 0,631276 \\
Random Forest & 0,054366 & 0,088403 & 0,790941 \\
Random Forest progresivo & 0,054337 & 0,088533 & 0,790326 \\
XGBoost & 0,054604 & 0,091735 & 0,774885 \\
XGBoost progresivo & 0,057295 & 0,113606 & 0,654747 \\
\bottomrule
\end{tabular}
\caption{Metricas finales de test por familia de modelo.}
\label{tab:resultados-modelos}
\end{table}

La regresion lineal presenta el peor RMSE y el menor $R^2$. Este resultado es coherente con el marco teorico: una superficie de volatilidad real tiene curvaturas, interacciones y comportamientos por regiones que dificilmente se capturan con una aproximacion lineal global. La linealidad sigue siendo util como control de complejidad y como referencia interpretable.

Random Forest y XGBoost muestran resultados muy similares en MAE, pero Random Forest reduce mas los errores grandes, como refleja el RMSE. El resultado sugiere que los patrones de volatilidad implicita del dataset se ajustan bien a particiones no lineales del espacio de features. La diferencia entre versiones progresivas y no progresivas no es uniforme: el progressive training mejora de forma importante la red tensor-train, pero no mejora Random Forest y empeora XGBoost en las ejecuciones disponibles.

\section{Interpretacion del rendimiento}

La diferencia entre MAE y RMSE es informativa. Varios modelos tienen MAE alrededor de 0,054--0,057, pero RMSE mas disperso. Esto indica que el error tipico puede ser parecido mientras que los errores extremos difieren. En finanzas de derivados, esos extremos importan: suelen concentrarse en regiones de baja liquidez, vencimientos muy cortos, moneyness extrema o episodios de mercado tensionado.

El $R^2$ de Random Forest y XGBoost confirma que los modelos capturan una parte sustancial de la variabilidad de test. No obstante, un $R^2$ cercano a 0,79 no significa que el modelo sea perfecto ni que pueda sustituir procedimientos de mercado. Significa que, dadas las variables utilizadas, el modelo explica gran parte de la variacion observada en la volatilidad implicita calculada. Quedan errores que deben analizarse con diagnosticos por region.

\section{Discusion sobre progressive training}

El progressive training parte de una hipotesis razonable: aprender primero las zonas mas cercanas a ATM puede estabilizar el modelo antes de incorporar alas. Los resultados muestran que esta idea no es universal. En la red tensor-train progresiva, el RMSE baja de 0,143025 a 0,117404 y el $R^2$ sube de 0,452787 a 0,631276. En cambio, Random Forest apenas cambia y XGBoost empeora frente a su version no progresiva.

La interpretacion es que la utilidad del entrenamiento progresivo depende de la dinamica interna de la familia. En redes, continuar entrenamiento por fases puede ordenar el aprendizaje de interacciones. En ensembles de arboles, la particion del espacio y los pesos por segmento ya capturan mucha heterogeneidad sin necesidad de fases. En boosting, repartir estimadores por fases puede limitar la capacidad de corregir residuos globales si la secuencia no queda bien equilibrada.

\section{Explicabilidad global}

La explicabilidad global responde a que variables sostienen las predicciones del modelo. El proyecto calcula SHAP mediante permutacion para mantener una interfaz comun entre familias. Esto permite comparar regresion lineal, Random Forest, XGBoost y redes bajo una misma semantica de contribucion aditiva.

En un modelo de volatilidad implicita, se espera que moneyness, vencimiento y sus interacciones tengan peso relevante. Si SHAP indicara dependencia dominante de variables sin interpretacion financiera o patrones incompatibles con la teoria, seria una senal de alerta. Por eso el dashboard no muestra solo importancias medias; tambien incluye dependence plots, heatmaps de contribuciones y comparacion visual entre variables.

Los arboles surrogate complementan SHAP al traducir el modelo principal a reglas. Una profundidad baja puede revelar cortes principales de moneyness y vencimiento. Si la fidelidad es baja, esas reglas solo sirven como intuicion parcial. Si una profundidad moderada logra buen RMSE de fidelidad, el comportamiento del modelo es mas resumible y, por tanto, mas comunicable.

La regresion simbolica busca una ecuacion cerrada que aproxime predicciones del modelo principal. Su valor esta en producir una expresion compacta para discusion, no en reemplazar automaticamente el modelo operativo. El dashboard conserva complejidad, perdida, score y metricas de fidelidad para evitar sobreinterpretar una formula atractiva pero imprecisa.

\section{Explicabilidad local}

La explicabilidad local permite analizar una prediccion concreta. El waterfall SHAP descompone la diferencia entre valor base y prediccion final. En opciones, esta lectura es util para explicar si una volatilidad alta se debe a moneyness extrema, vencimiento corto, tipo de opcion o combinaciones de variables.

Los vecinos cercanos aportan soporte empirico. Una prediccion puede tener contribuciones SHAP aparentemente razonables y aun asi estar lejos de observaciones historicas. El proyecto calcula vecinos contra el split de train transformado para evitar usar el mismo test como soporte. La proyeccion PCA ayuda a visualizar si la muestra cae en una region poblada o aislada.

Esta combinacion reduce un riesgo frecuente de la explicabilidad local: explicar matematicamente cualquier salida aunque el punto este fuera de distribucion. En un sistema financiero, una prediccion fuera de soporte debe marcarse como menos fiable, aunque el modelo devuelva una cifra.

\section{Superficies, ICE y ALE}

La pestana Behaviour and Surface estudia como cambia la volatilidad predicha al mover moneyness, vencimiento, strike, subyacente o tipo. Las superficies de volatilidad se generan alrededor de anclas reales del test. Esto evita construir superficies completamente abstractas y mantiene el analisis ligado a condiciones observadas.

Los heatmaps y superficies 3D muestran forma global local: smiles, pendientes y cambios por vencimiento. Las curvas smile resumen cortes por moneyness; las term structures resumen cortes por vencimiento. Los checks financieros avisan de posibles incoherencias, como patrones excesivamente irregulares.

ICE y ALE cumplen funciones distintas. ICE muestra respuestas individuales y permite observar heterogeneidad entre contratos. ALE resume efectos medios locales reduciendo extrapolacion en regiones con poca densidad. En conjunto, ambas herramientas ayudan a distinguir una sensibilidad general de comportamientos particulares.

\section{Diagnostico de errores}

La pestana Diagnosis presenta metricas agregadas, scatter predicted vs actual, heatmaps de residuos y errores por moneyness y vencimiento. Esta capa es necesaria porque una metrica global puede ocultar fallos sistematicos. Un modelo con buen RMSE medio puede subestimar volatilidad en alas, sobreestimar en vencimientos cortos o degradarse en episodios concretos.

El residual se define como $y_i-\hat{y}_i$. Residuos positivos indican subestimacion de volatilidad; residuos negativos, sobreestimacion. El error absoluto mide magnitud local. Los mapas por moneyness y vencimiento permiten identificar regiones donde el modelo requiere mas datos, mejores features o restricciones financieras adicionales.

\section{Respuesta a la pregunta de investigacion}

La pregunta inicial era si se podia construir una herramienta que estimara volatilidad implicita y explicara sus predicciones. Los resultados disponibles apoyan una respuesta afirmativa con matices. La estimacion es razonable fuera de muestra para familias no lineales, especialmente Random Forest y XGBoost. La explicabilidad esta integrada en el flujo y cubre niveles global, local, funcional y de diagnostico.

El matiz es que la calidad depende del dominio de datos y de la region de la superficie. El sistema no convierte automaticamente una prediccion en verdad de mercado. La volatilidad implicita de entrada se calcula con Black-76 y filtros especificos; el modelo aprende esa construccion. Por tanto, la interpretacion debe entenderse como explicacion de un modelo entrenado sobre una base construida, no como descubrimiento causal de precios.

\section{Limitaciones}

La primera limitacion es la dependencia de Black-76. El modelo de inversion es adecuado para opciones sobre futuros, pero impone supuestos. Las volatilidades implicitas calculadas heredan esos supuestos y cualquier error de inputs.

La segunda limitacion es la microestructura. Aunque se filtra por lag de subyacente y tipo de operacion, las operaciones reales pueden reflejar liquidez desigual, horquillas, episodios de estres o negociaciones no representativas. Algunas volatilidades extremas probablemente proceden de estas condiciones.

La tercera limitacion es la ausencia de variables adicionales de mercado. El modelo usa tipo de opcion, strike, subyacente, vencimiento y tipo. Podrian incorporarse variables de liquidez, hora de sesion, volumen, volatilidad realizada, indicadores macro o informacion de order book si estuvieran disponibles.

La cuarta limitacion es la interpretacion de los explicadores. SHAP, ICE, ALE, surrogates y regresion simbolica explican el comportamiento del modelo bajo ciertos supuestos y muestras. No prueban causalidad ni garantizan estabilidad fuera del dominio observado.

\section{Comparacion detallada de modelos}

La regresion lineal obtiene RMSE de test 0,183824. Su $R^2$ de 0,096059 indica que apenas mejora a un predictor constante en el bloque final. La version progresiva empeora ligeramente. Esto sugiere que segmentar el entrenamiento no soluciona la rigidez estructural del modelo. Aunque las features incluyen no linealidades basicas, la relacion global sigue siendo demasiado compleja para una combinacion lineal.

La red tensor-train mejora claramente la regresion lineal. La version no progresiva alcanza RMSE 0,143025 y $R^2$ 0,452787; la progresiva alcanza RMSE 0,117404 y $R^2$ 0,631276. Esta mejora muestra que ordenar el aprendizaje por segmentos de moneyness puede ayudar a una arquitectura neuronal. Aun asi, el resultado queda lejos de los ensembles de arboles, probablemente porque los datos tabulares y el tamano de muestra favorecen particiones no lineales directas.

Random Forest presenta la mejor metrica final. La version progresiva tiene MAE levemente menor, 0,054337 frente a 0,054366, pero RMSE levemente mayor, 0,088533 frente a 0,088403. La diferencia es pequena y no justifica preferir la variante progresiva si se prioriza RMSE. El modelo no progresivo queda como mejor candidato por simplicidad y rendimiento.

XGBoost logra RMSE 0,091735, muy cercano a Random Forest. Su version progresiva cae a RMSE 0,113606. La diferencia sugiere que el boosting se beneficia de optimizar globalmente la secuencia de arboles, mientras que dividir estimadores por segmentos puede limitar su capacidad de corregir errores residuales en todo el espacio.

\section{Generalizacion temporal}

El test cubre desde septiembre de 2021 hasta junio de 2022, periodo posterior al trainval. Esta separacion es exigente porque evalua el modelo en fechas no usadas durante seleccion. La diferencia entre trainval y test muestra degradacion esperable. Por ejemplo, Random Forest pasa de RMSE trainval 0,070600 a test 0,088403; XGBoost de 0,074342 a 0,091735. La brecha existe, pero no destruye el rendimiento.

En la regresion lineal, la brecha es mayor en terminos relativos y el $R^2$ test queda muy bajo. Esto refuerza la idea de que un modelo con sesgo alto no solo falla en entrenamiento, sino que tampoco generaliza bien ante cambios temporales. En redes, la brecha tambien es relevante, lo que puede indicar sensibilidad a regimenes o necesidad de regularizacion y variables adicionales.

\section{Lectura financiera de los resultados}

El mejor comportamiento de Random Forest y XGBoost es coherente con una superficie por regiones. En volatilidad implicita, la sensibilidad puede cambiar bruscamente entre opciones ATM y alas, vencimientos cortos y medios, calls y puts. Los arboles capturan este tipo de particiones de forma natural. Ademas, no requieren que las variables esten escaladas ni que la relacion sea suave globalmente.

Sin embargo, que un modelo de arboles funcione bien no implica que la superficie real sea discontinua. La prediccion final de un ensemble puede aproximar funciones suaves mediante muchas particiones. Por eso las superficies generadas en el dashboard son necesarias: permiten observar si el modelo aprendido presenta saltos o irregularidades no deseadas.

La volatilidad implicita media de la base es 0,2061, con mediana 0,1738. Un RMSE de 0,0884 representa una fraccion importante de la volatilidad media, pero debe interpretarse considerando que la distribucion tiene cola larga y maximos extremos. El MAE de 0,0544 indica que el error tipico absoluto es menor que el RMSE y que parte de la penalizacion procede de errores grandes.

\section{Discusion sobre datos y filtros}

La reduccion desde 428.448 operaciones de opciones con subyacente hasta 312.615 volatilidades finales muestra que el solver y los filtros eliminan una parte relevante. Esta reduccion es esperable: no toda operacion produce una volatilidad implicita estable dentro del rango definido, y no toda asociacion de subyacente cumple el filtro final.

El filtro de lag maximo de 180 minutos es conservador en comparacion con la mediana final de 0,080 minutos. La media de 7,991 minutos indica que algunas observaciones tienen lags mas altos, pero la mayor parte de la base usa subyacentes muy cercanos temporalmente. Esto mejora la credibilidad de la inversion Black-76, porque el precio del futuro usado refleja condiciones proximas a la operacion de opcion.

El tiempo hasta vencimiento medio de 23,44 dias indica predominio de vencimientos relativamente cortos, coherente con opciones mensuales negociadas. El maximo de 623 dias en la base final revela que tambien hay observaciones largas, aunque no dominan. Esta mezcla justifica incluir vencimiento y raiz de vencimiento como variables, y diagnosticar errores por madurez.

\section{Coherencia de explicaciones}

Un resultado cuantitativo aceptable debe ir acompanado de explicaciones coherentes. En un modelo financiero, seria sospechoso que una variable irrelevante dominara SHAP o que la respuesta a moneyness fuera contraria a patrones esperados sin razon de mercado. La documentacion del dashboard establece precisamente esa lectura: usar visualizaciones SHAP para detectar drivers, ICE/ALE para respuesta y superficies para coherencia financiera.

La existencia de modelos equivalentes aporta una segunda prueba. Si un arbol surrogate de baja profundidad logra fidelidad razonable, se puede comunicar el comportamiento principal con reglas sencillas. Si necesita profundidad alta, el modelo es mas complejo y conviene evitar narrativas simplificadas. La regresion simbolica cumple un papel similar, pero con formulas en lugar de reglas.

\section{Uso ante un tribunal o usuario profesional}

La memoria puede defenderse mostrando tres evidencias. Primero, la cadena de datos esta materializada y validada: no se parte de una tabla opaca. Segundo, la evaluacion temporal reserva test final y controla contratos compartidos. Tercero, la explicabilidad no se limita a importancias, sino que integra diagnostico de rendimiento, comportamiento de superficie y soporte local.

Ante un usuario profesional, la recomendacion seria no presentar el modelo como motor autonomo de precios. Es una herramienta de analisis y explicacion de volatilidad implicita calculada. Puede servir para identificar patrones, comparar familias, estudiar regiones de error y generar predicciones orientativas. Para uso de produccion harian falta validaciones adicionales, monitorizacion continua y controles de calidad intradia.

\section{Riesgos residuales}

Persisten riesgos incluso con el protocolo aplicado. Puede existir cambio de regimen posterior a junio de 2022. Puede haber sesgos en las operaciones disponibles frente al universo total de cotizaciones. Puede haber dependencia de convenciones de vencimiento o de tipos. Puede haber regiones con pocas observaciones donde el modelo interpola o extrapola sin suficiente soporte.

La respuesta del proyecto a estos riesgos no es eliminarlos por completo, sino hacerlos visibles. Los vecinos muestran soporte local; los heatmaps muestran errores por region; las superficies muestran irregularidades; los metadatos conservan configuracion. Esta filosofia es adecuada para un TFM aplicado: construir una herramienta que no solo predice, sino que tambien expone sus propias zonas de duda.

\section{Escenarios de uso}

Un primer escenario de uso es la revision de un modelo candidato. El usuario selecciona una familia, revisa metricas finales, observa predicted vs actual y despues inspecciona SHAP global. Si las variables dominantes son coherentes con teoria financiera, avanza a superficies y diagnosticos. Si aparecen dependencias inesperadas, vuelve a revisar features o datos.

Un segundo escenario es el analisis de una operacion concreta. El usuario selecciona una muestra de test, observa volatilidad real, predicha y residual. El waterfall SHAP indica que variables empujaron la prediccion. Los vecinos muestran si existen operaciones similares en entrenamiento. Si el residual es alto y los vecinos son escasos, la conclusion no es solo que el modelo fallo, sino que la region puede estar mal representada.

Un tercer escenario es la comunicacion de comportamiento global. Los arboles surrogate y la regresion simbolica permiten presentar reglas o formulas aproximadas. Esta comunicacion es util en defensa academica y en reuniones de validacion, siempre que se acompane de fidelidad. Una regla simple con baja fidelidad puede ser pedagogica, pero no debe usarse para justificar decisiones operativas.

\section{Comparacion con un enfoque puramente teorico}

Un enfoque puramente teorico podria calibrar una superficie parametrica o un modelo de volatilidad local. Ese enfoque tiene ventajas: puede incorporar restricciones financieras y producir superficies suaves. Sin embargo, exige decisiones de parametrizacion y calibracion. El enfoque de este TFM es distinto: aprende una aproximacion supervisada desde datos de operaciones y despues audita su comportamiento.

Ambos enfoques no son excluyentes. Un trabajo futuro podria usar los modelos de machine learning como herramienta de interpolacion inicial y despues imponer restricciones de ausencia de arbitraje. Tambien podria usar modelos teoricos para generar features adicionales o benchmarks. La contribucion de este TFM es mostrar que la aproximacion data-driven puede ser trazable y explicable.

\section{Discusion sobre ausencia de resultados visuales en la memoria}

El repositorio contiene un dashboard diseñado para visualizaciones interactivas. Esta memoria prioriza tablas, descripciones y anexos porque se construye a partir de la documentacion disponible y de metadatos locales. En una version final compilada seria recomendable exportar capturas de las pestanas principales: SHAP global, superficie de volatilidad, waterfall local, vecinos y diagnostico.

La ausencia de capturas no impide documentar el sistema, pero limita la comunicacion visual. Las figuras permitirian comprobar rapidamente si las superficies son suaves, si los residuos se concentran en regiones concretas y si las contribuciones SHAP son financieramente razonables. Por tanto, se recomienda anadirlas cuando se estabilice el modelo final elegido.

\section{Sintesis interpretativa}

La sintesis de resultados es la siguiente. La cadena de datos produce una muestra amplia y temporalmente ordenada. Los modelos lineales no capturan suficiente estructura. Los ensembles de arboles son los mejores candidatos. Las redes muestran potencial, especialmente con entrenamiento progresivo, pero no lideran las metricas. La explicabilidad es suficientemente rica para diagnosticar comportamiento y comunicar limitaciones.

El resultado no debe presentarse como cierre definitivo. Es una base solida para iterar. La mejora mas prometedora no necesariamente sera cambiar de algoritmo, sino enriquecer variables y mejorar controles de superficie. Con datos de liquidez y capturas visuales del dashboard, el trabajo podria evolucionar hacia una herramienta mas cercana a validacion profesional.
